\chapter{Introduction}
\label{ch:introduction}

\acrfull{y90} \acrfull{sirt}, also referred to as \acrfull{tare}, is an established treatment for primary and secondary liver malignancies \cite{sundramSelectiveInternalRadiation2017}. In this procedure, \gls{y90}-labelled microspheres are administered into the hepatic artery to preferentially irradiate tumour tissue while sparing healthy liver parenchyma. The administered activities are high (typically of the order \SI{2}{\giga\becquerel}), and clinically relevant absorbed-dose heterogeneity can arise at multiple spatial scales due to vascular delivery, microsphere clustering, and patient-specific hepatic anatomy \cite{strigariEvidenceBaseUse2014}. Tumour doses of clinical interest often exceed \SI{200}{\gray} \cite{salemClinicalDosimetricConsiderations2019}. Further treatment can increase the dose to healthy liver tissue, making treatment planning and verification important \cite{reinckeHepaticDecompensationTransarterial2022}. Post-therapy quantitative imaging and dosimetry support treatment verification, outcome modelling, and the planning of subsequent therapy.

Despite its importance, routine post-\gls{sirt} dosimetry is not universally performed in clinical practice, primarily due to workflow constraints and limitations in quantitative accuracy \cite{kappadathQuantitativeEvaluation90YPET2024, gearEANMEnablingGuide2023}. Post-therapy imaging of \gls{y90} is limited by low effective count statistics and less favourable imaging physics than standard diagnostic radionuclides \cite{kappadathQuantitativeEvaluation90YPET2024}. These compound the statistical noise, system-modelling errors and hardware limitations encountered in emission tomography \cite{baiTumorQuantificationClinical2013}. Improving reconstruction in this setting could support more reliable activity and absorbed-dose estimates \cite{davisPersonalisationMolecularRadiotherapy2020}.

In post-\gls{sirt} imaging, the two clinically relevant emission modalities are complementary. \gls{y90} \acrfull{pet} offers comparatively high spatial resolution, but is count-limited because of the low internal pair-production branch. Bremsstrahlung \acrfull{spect}, by contrast, provides higher count statistics but suffers from poor spatial resolution and modelling challenges, particularly for scatter and septal penetration \cite{kappadathQuantitativeEvaluation90YPET2024}. These emission images are routinely paired with \acrfull{ct}, which provides attenuation correction and anatomical context. This motivates reconstruction approaches in which information is shared across modalities, with \glsentryshort{ct} contributing structural guidance, and \glsentryshort{pet} and \glsentryshort{spect} contributing complementary functional information. Synergistic methods, where such information is incorporated during reconstruction rather than only after image formation, have been investigated in related hybrid imaging contexts \cite{arridgeOverviewSynergisticReconstruction2021}. For post-\gls{y90} \gls{sirt} imaging, the question is whether higher-count but blurred \gls{spect} data can support a low-count \gls{pet} reconstruction when combined with \gls{ct}-derived directional guidance.

This thesis develops and evaluates a synergistic reconstruction framework for post-\gls{y90} \gls{sirt}, examining whether functional coupling and anatomical directional guidance together improve reconstruction beyond either mechanism alone. To make bremsstrahlung \gls{spect} useful in this setting, the work combines established scatter-estimation methods with refinements to their application and an effective range-blur model. It introduces a weighted directional total nuclear variation prior for joint \glsentryshort{pet}/\glsentryshort{spect} reconstruction with \glsentryshort{ct}-derived directional guidance, and studies subset scheduling and prior-aware preconditioning for the resulting penalised-likelihood problem. Evaluation focuses on image recovery, lesion contrast and noise; improved dosimetry is a motivation rather than an outcome demonstrated here.

The thesis is organised as follows:
\begin{itemize}
    \item \textbf{Chapter 1} motivates synergistic post-\gls{sirt} imaging and outlines the thesis contributions.
    \item \textbf{Chapter 2} introduces the physics and measurement models of \glsentryshort{pet}, \glsentryshort{spect}, and \glsentryshort{ct}, with emphasis on post-\gls{y90} imaging.
    \item \textbf{Chapter 3} formulates penalised-likelihood emission reconstruction and reviews reconstruction priors and optimisation algorithms.
    \item \textbf{Chapter 4} develops practical improvements to the \gls{y90} bremsstrahlung \glsentryshort{spect} system model.
    \item \textbf{Chapter 5} studies anatomy-informed regularisation in a controlled post-reconstruction deconvolution setting.
    \item \textbf{Chapter 6} introduces the weighted directional total nuclear variation prior and its differentiable formulation.
    \item \textbf{Chapter 7} investigates the optimisation of \glsentryshort{wdtnv} through controlled single-bed phantom experiments on subset organisation and prior-aware preconditioning.
    \item \textbf{Chapter 8} extends the framework to multi-bed patient acquisitions and evaluates \glsentryshort{wdtnv} on phantom and clinical data for joint post-\gls{y90} \glsentryshort{pet}/\glsentryshort{spect} reconstruction with \glsentryshort{ctac}-derived guidance.
    \item \textbf{Chapter 9} discusses the findings, limitations, and directions for future work, and concludes the thesis.
    \item \textbf{The Appendix} documents additional software and reproducibility material that supports the thesis workflow.
\end{itemize}
